\section{Strict positivity between rational pencil parameters}

\begin{lemma}[Derivative tests on a dense rational set]
\label{lem:nef-ample-polynomial-rational-derivative-nonnegative}
\lean{EnvelopeProperty.SeshadriAmplenessCriterionNefAmple.polynomial_derivative_nonneg_of_positive_rational_tests}
\leanok
Let $g\in\mathbb R[T]$. If $g'(q)\geq0$ for every positive rational
$q$, then $g'(t)\geq0$ for every positive real $t$.
\end{lemma}
\begin{proof}
If $g'(t)<0$ for a positive real $t$, continuity of the polynomial
$g'$ gives an open neighborhood of $t$ consisting of positive points
where $g'$ is negative. This neighborhood contains a rational number,
contradicting the stated rational tests.
\end{proof}

\begin{theorem}[Strict positivity at every positive real parameter]
\label{thm:nef-ample-polynomial-real-positivity}
\lean{EnvelopeProperty.SeshadriAmplenessCriterionNefAmple.polynomial_pos_of_positive_rational_tests_and_derivative_nonneg}
\leanok
\uses{lem:nef-ample-polynomial-rational-derivative-nonnegative}
Let $g\in\mathbb R[T]$. Suppose that $g(q)>0$ and $g'(q)\geq0$
for every positive rational $q$. Then $g(t)>0$ for every positive
real $t$.
\end{theorem}
\begin{proof}
The preceding lemma makes $g'$ nonnegative throughout $(0,\infty)$.
The mean value theorem therefore makes $g$ nondecreasing on this interval.
For any $t>0$, choose a rational $q$ with $0<q<t$. The hypotheses and
monotonicity give $0<g(q)\leq g(t)$, as required. This step uses the
derivative hypothesis: positivity on the rational points by itself would
not exclude a zero at an irrational parameter.
\end{proof}
